% ============================================================
% Manuscript: Towards Coloc-Confirmed TWAS of Ankylosing Spondylitis
%             to Prioritize Tissue-Specific Drug Targets
%
% Target journals (adapt formatting as needed):
%   - Human Genomics (BioMed Central)
%   - BMC Bioinformatics (BioMed Central)
%   - Frontiers in Genetics (Computational Genomics section)
%
% NOTE TO AUTHOR:
%   Sections marked [PLACEHOLDER] contain draft text and must be
%   updated with real results after pipeline analyses are complete.
%   Do NOT submit manuscript text from placeholders as scientific findings.
% ============================================================

\documentclass[12pt,a4paper]{article}

% ── Packages ─────────────────────────────────────────────────
\usepackage[margin=2.5cm]{geometry}
\usepackage{amsmath}
\usepackage{graphicx}
\usepackage{booktabs}
\usepackage{hyperref}
\usepackage{natbib}
\usepackage{setspace}
\usepackage{lineno}
\usepackage{xcolor}

% Highlight placeholders in draft
\newcommand{\placeholder}[1]{\textcolor{red}{\textbf{[PLACEHOLDER: #1]}}}

\doublespacing
\linenumbers

% ── Title and authors ─────────────────────────────────────────
\title{Towards Colocalisation-Confirmed Transcriptome-Wide Association Study
       of Ankylosing Spondylitis to Prioritize Tissue-Specific Drug Targets}

\author{
      Bhargav Yellepeddi$^{1}$ \\
  \\
      $^{1}$Department and institution details to be finalized before submission \\
  \\
      Correspondence: to be finalized before submission
}

\date{\today}

\begin{document}

\maketitle

% ── Abstract ─────────────────────────────────────────────────
\begin{abstract}
\textbf{Background:}
Ankylosing spondylitis (AS) is a chronic inflammatory spondyloarthropathy
with substantial genetic heritability and growing GWAS-identified risk loci.
Despite this, the causal genes and effector tissues underlying AS risk remain
largely uncharacterized.

\textbf{Methods:}
We integrated AS GWAS summary statistics with GTEx v8 gene expression
prediction models across 8 tissues using S-PrediXcan.
Global Benjamini-Hochberg false-discovery-rate correction was applied across
all gene-tissue pairs. For FDR-significant hits, we performed Bayesian
colocalization (coloc.abf) to distinguish genuinely co-localised GWAS and
eQTL signals (PP4) from coincidental proximity (PP3).

\textbf{Results:}
S-PrediXcan yielded 2,313 valid gene-tissue tests across the 8 configured
GTEx v8 tissues. After global FDR correction, 248 gene-tissue associations
were retained for follow-up. Bayesian colocalization identified 35
gene-tissue pairs spanning 18 unique genes with PP4 $\geq 0.7$, with Whole
Blood accounting for 11 of the 35 prioritized pairs. Top colocalized signals
included \textit{ERAP1}, \textit{GPR25}, \textit{CARD9}, and \textit{FAS}.

\textbf{Conclusions:}
Combining TWAS with Bayesian colocalization reduced a broad set of
transcriptome associations to a smaller, higher-confidence set of
tissue-specific candidate genes for AS follow-up studies.

\textbf{Keywords:}
ankylosing spondylitis; TWAS; S-PrediXcan; colocalization; coloc; eQTL; GTEx; drug target prioritization
\end{abstract}

\newpage

% ── Introduction ──────────────────────────────────────────────
\section{Introduction}

Ankylosing spondylitis (AS) is a seronegative inflammatory arthritis
predominantly affecting the axial skeleton, with a population prevalence of
approximately 0.1--0.5\% in European populations \citep{Reveille2011}.
The disease exhibits high heritability (estimated 70--90\%), with HLA-B*27
as the largest genetic risk factor, accounting for $\sim$20--30\% of
genetic variance \citep{Brown2010}.
Genome-wide association studies (GWAS) have identified over 100 risk loci
beyond the MHC region, yet the majority of variants lie in non-coding regions,
complicating interpretation of their molecular mechanisms \citep{Ellinghaus2016}.

Transcriptome-wide association studies (TWAS) offer a complementary approach
to GWAS by integrating expression quantitative trait loci (eQTLs) with
GWAS summary statistics to identify genes whose predicted expression is
associated with disease risk \citep{Gamazon2015,Barbeira2018}.
S-PrediXcan (MetaXcan) implements a summary-statistic-based approach that
enables multi-tissue TWAS without requiring individual-level genotype data
\citep{Barbeira2018}.
The GTEx v8 resource provides high-quality eQTL prediction models trained
on 49 tissues from 838 post-mortem donors, enabling tissue-specific
prioritization of disease genes \citep{GTExConsortium2020}.

A critical limitation of TWAS is that statistical associations may arise from
coincidental linkage disequilibrium rather than a shared underlying causal
variant.
Bayesian colocalization using coloc.abf \citep{Giambartolomei2014} formally
tests whether the GWAS signal and the eQTL signal at a given locus are driven
by the same variant (posterior probability of hypothesis H4; PP4).
High PP4 (typically $\geq 0.7$) provides strong evidence that the disease
GWAS signal and the tissue-specific eQTL signal co-localize, supporting the
gene as a causal driver of AS risk in that tissue.

Here, we describe a reproducible pipeline integrating GWAS summary statistics
for AS with GTEx v8 S-PrediXcan TWAS followed by coloc-based prioritization,
with the aim of identifying tissue-specific candidate drug targets for AS.
The pipeline is designed for transparency, modularity, and extensibility to
future AS GWAS or eQTL data.

% ── Methods ──────────────────────────────────────────────────
\section{Methods}

\subsection{GWAS Summary Statistics}

Cohort-level genotyping, imputation-panel, and study-specific QC details
should be reported directly from the source GWAS publication rather than
inferred from downstream pipeline artifacts.

We obtained AS GWAS summary statistics from the NHGRI-EBI GWAS Catalog
accession GCST005529 \citep{Ellinghaus2016}. The dataset comprised 10,619
cases and 15,145 controls (25,764 total samples). Summary statistics were
originally reported on genome build hg19.

\subsection{GWAS Harmonization}

Raw summary statistics were harmonized to hg38 using the following steps:
\begin{enumerate}
  \item Column names were standardized to a canonical schema
        (varID, chrom, pos, effect\_allele, non\_effect\_allele, beta, se, pvalue, N).
  \item Where applicable, coordinates were lifted from hg19 to hg38 using
        \texttt{pyliftover}.
  \item GTEx-style variant identifiers (varIDs) were constructed as
        \texttt{chr\{chrom\}\_\{pos\}\_\{ref\}\_\{alt\}\_b38}.
        This format is required for matching GWAS variants to GTEx
        PredictDB prediction models.
  \item Ambiguous palindromic single-nucleotide polymorphisms (AT/CG) were
        removed by default to prevent misalignment of effect directions across
        strands, as these cannot be reliably strand-disambiguated without
        allele frequency information.
  \item Variants with missing effect sizes, standard errors, or p-values were
        excluded.
\end{enumerate}

\subsection{Transcriptome-Wide Association Study}

TWAS was performed using S-PrediXcan (MetaXcan repository commit
\texttt{9b9f76b})
\citep{Barbeira2018} with GTEx v8 mashr-based prediction models
\citep{GTExConsortium2020,Barbeira2019}.
Eight tissues relevant to AS pathophysiology were analysed: Whole Blood,
Adipose Subcutaneous, Muscle Skeletal, Colon Sigmoid, Colon Transverse,
Small Intestine Terminal Ileum, Spleen, and Lung.

To control for multiple testing, Benjamini-Hochberg false-discovery rate
(FDR) correction was applied globally across all gene-tissue pairs tested.
Gene-tissue pairs with FDR-adjusted q-value $< 0.05$ were retained as
TWAS-significant hits.

\subsection{Bayesian Colocalization}

For each TWAS-significant gene-tissue pair, we performed Bayesian
colocalization using the coloc.abf function from the \texttt{coloc} R
package (v5.2.3) \citep{Giambartolomei2014}.
Regional GWAS summary statistics ($\pm$1\,Mb around the gene transcription
start site) and tissue-matched GTEx v8 eQTL summary statistics were extracted.
All SNPs within the region were included without pre-filtering by p-value,
as removing non-significant SNPs inflates PP4 estimates \citep{Wallace2021}.

Effect sizes (beta) and their variances ($\sigma^2 = \mathrm{SE}^2$;
\textit{varbeta}) were provided as inputs to coloc.abf.
Standard coloc priors were used: $p_1 = 10^{-4}$, $p_2 = 10^{-4}$,
$p_{12} = 10^{-5}$.

Coloc results were interpreted using three complementary metrics:
\begin{itemize}
  \item PP4 $\geq 0.7$: primary evidence for a shared causal variant.
  \item PP3 + PP4 $\geq 0.8$: powered locus with sufficient signal to
        distinguish hypotheses H3 and H4.
  \item PP4 / (PP3 + PP4) $\geq 0.9$: strong preference for shared
        signal over distinct signals.
\end{itemize}

For loci with potential multiple independent signals, coloc.susie was
reserved as planned follow-up sensitivity analysis rather than part of the
current reported run \citep{Wallace2021}.

\subsection{Reproducibility}

All analyses were performed using a custom Python pipeline
(available at \url{https://github.com/yellepeddibk/as-twas-coloc}).
Each run saves a provenance manifest containing the pipeline version,
git commit hash, configuration snapshot, coloc priors, and input file checksums.
The current implementation also records Python package versions together with
external tool identifiers, including the MetaXcan repository commit, R
version, and coloc package version when available.
Random seeds were fixed throughout.
The pipeline runs in a lightweight mock mode for CI/CD validation without
requiring external data.

% ── Results ──────────────────────────────────────────────────
\section{Results}

% Numerical values below are limited to counts validated against current
% result artifacts in this repository.

\subsection{GWAS Harmonization and QC}

The configured workflow ingested GWAS Catalog accession GCST005529 on hg19
coordinates, lifted variants to hg38, standardized the summary-statistics
schema, removed ambiguous palindromic SNPs by default, and generated a
GTEx-compatible varID-indexed file for downstream TWAS. Final per-step SNP
accounting should be reported from the production run artifacts alongside the
submitted manuscript.

\subsection{TWAS Identifies Association Signals Across Multiple Tissues}

Across the eight configured GTEx tissues, S-PrediXcan returned 2,316 raw
gene-tissue rows. Three rows had missing p-values and were excluded before
multiple-testing correction, leaving 2,313 valid gene-tissue tests. Global
Benjamini-Hochberg correction identified 248 significant associations
($q < 0.05$), which were carried forward to colocalization. These signals
spanned multiple chromosomes and tissues, with the strongest downstream
colocalization support concentrated in Whole Blood.

% Insert TWAS QQ-plot
% \begin{figure}[htbp]
%   \centering
%   \includegraphics[width=0.6\textwidth]{../../results/figures/qqplot_twas.png}
%   \caption{QQ-plot of TWAS p-values across all gene-tissue pairs.}
%   \label{fig:qqplot}
% \end{figure}

% Insert Manhattan-style plot
% \begin{figure}[htbp]
%   \centering
%   \includegraphics[width=\textwidth]{../../results/figures/manhattan_like_twas.png}
%   \caption{Manhattan-style plot of TWAS associations.}
%   \label{fig:manhattan}
% \end{figure}

\subsection{Colocalization Prioritizes a Subset of Genes with Shared Signals}

All 248 TWAS-significant gene-tissue pairs underwent colocalization analysis.
Thirty-five gene-tissue pairs, representing 18 unique genes, met the primary
shared-signal threshold of PP4 $\geq 0.7$. Whole Blood accounted for 11 of the
35 prioritized pairs. The highest-confidence signals included
\textit{ERAP1} in Whole Blood (PP4 = 0.999) and Small Intestine Terminal
Ileum (PP4 = 0.999), \textit{GPR25} in Whole Blood (PP4 = 0.997),
\textit{CARD9} in Whole Blood (PP4 = 0.974), and \textit{FAS} in Colon
Sigmoid (PP4 = 0.973).

% Insert coloc summary barplot
% \begin{figure}[htbp]
%   \centering
%   \includegraphics[width=0.7\textwidth]{../../results/figures/coloc_summary.png}
%   \caption{Top coloc-confirmed gene-tissue pairs ranked by PP4.}
%   \label{fig:coloc_barplot}
% \end{figure}

% Insert top-genes table
% \input{top_genes_table.tex}  % generated by pipeline via manuscript_support.py

% ── Discussion ───────────────────────────────────────────────
\section{Discussion}

Several prioritized genes align with established AS biology while also
highlighting follow-up candidates. \textit{ERAP1}, one of the best-supported
AS loci, remained the top colocalized signal in both Whole Blood and Small
Intestine Terminal Ileum, which supports the pipeline's ability to recover
known disease biology. Additional prioritized genes such as \textit{CARD9},
\textit{UBE2L3}, and \textit{FAS} implicate innate immune signaling,
ubiquitin-mediated regulation, and apoptosis-related pathways with plausible
roles in chronic inflammation. Whole Blood contributed the largest number of
coloc-confirmed pairs, consistent with a systemic immune component, while the
presence of intestine- and colon-associated signals argues against restricting
interpretation to peripheral blood alone. Less-characterized candidates such
as \textit{GPR25}, \textit{CCDC116}, and \textit{NPIPB7} are therefore best
treated as prioritized hypotheses for replication, fine-mapping, and
functional follow-up rather than definitive causal assignments.

% ── Limitations ──────────────────────────────────────────────
\section{Limitations}

The current analysis has several limitations that should be considered when
interpreting findings.
First, the GTEx v8 eQTL models are derived from post-mortem tissues, which
may not fully capture in vivo inflammatory states relevant to AS.
Second, the application of coloc.abf assumes a single causal variant per locus;
for loci with multiple independent signals, coloc.susie should be applied.
Third, the GTEx sample sizes for individual tissues may be insufficient to
detect eQTLs for all disease-relevant genes.
Fourth, three S-PrediXcan output rows had missing p-values and were excluded
from global FDR correction, underscoring the need to report valid tested-pair
counts directly rather than raw output row totals.

% ── Future Work ──────────────────────────────────────────────
\section{Future Work}

Planned next steps include multi-signal colocalization with coloc.susie,
fine-mapping of top loci, replication in independent AS cohorts,
integration of regulatory annotations such as chromatin accessibility, and
drug-target mapping for the highest-confidence prioritized genes.

% ── References ───────────────────────────────────────────────
\section*{Acknowledgements}

We thank Dr.\ Ratul Chowdhury for mentorship and project guidance. AS GWAS
summary statistics were obtained from the NHGRI-EBI GWAS Catalog
(GCST005529), and tissue expression models were derived from GTEx v8.
Funding and institutional support details should be finalized before
submission.

\bibliographystyle{unsrtnat}
\bibliography{references}

\end{document}
